\section{Introdução}
\label{sec:introducao}

\subsection{Contextualização}

A Engenharia Econômica para Software trata das decisões de investimento em
sistemas sob incerteza: quanto custa construir, em quanto tempo e se o retorno
esperado justifica o desembolso. Todas essas perguntas dependem de uma
estimativa de esforço, já que o esforço, valorado pelo custo do tempo da
equipe, é o principal componente do custo de um projeto de software
\cite{boehm1981}. Não surpreende que a estimativa de custos seja um dos temas
mais estudados da área: \citeonline{jorgensen2007} identificaram 304 artigos
sobre o assunto em 76 periódicos.

Nas últimas duas décadas, boa parte desse planejamento migrou de planos
detalhados e preditivos para ciclos curtos e adaptativos. Desde o Manifesto
Ágil \cite{agilemanifesto2001}, métodos como Scrum e Extreme Programming (XP)
se difundiram: o 17º \textit{State of Agile Report}, pesquisa com 788
profissionais, registra que 71\% dos respondentes usam práticas ágeis no ciclo
de desenvolvimento de software \cite{digitalai2024}. Nesses métodos, os
requisitos são escritos como histórias de usuário, técnica originária do XP
\cite{beck2004,cohn2004}, e estimados em \textit{story points}, unidades
relativas de tamanho e complexidade, e não em horas ou em reais
\cite{cohn2005}. A soma dos pontos concluídos em uma iteração define a
\textit{velocity} da equipe, que serve de base para prever quantas iterações
faltam para esvaziar o \textit{backlog} \cite{cohn2005}. O Guia do Scrum, por
sua vez, não prescreve uma técnica de dimensionamento e atribui essa
responsabilidade aos desenvolvedores \cite{scrumguide2020}.

\subsection{Problema e relevância econômica}

Como o \textit{story point} não é uma unidade monetária, é preciso uma ponte
entre a escala de pontos e a escala financeira. A ponte mais simples combina
três grandezas: o \textit{backlog} restante $B$ (em pontos), a \textit{velocity}
média $\bar{v}$ (pontos por \textit{sprint}) e o custo da equipe por
\textit{sprint} $c_s$. Com elas obtêm-se o número previsto de \textit{sprints}
$N$, o custo por ponto $c_p$ e o custo total previsto $C$:
%
\begin{equation}
  N = \frac{B}{\bar{v}}, \qquad
  c_p = \frac{c_s}{\bar{v}}, \qquad
  C = N\,c_s = B\,c_p .
  \label{eq:previsao}
\end{equation}
%
Por exemplo, com valores hipotéticos de $B = 240$ pontos,
$\bar{v} = 30$ pontos por \textit{sprint} e $c_s = \text{R\$}\,60.000$,
tem-se $N = 8$ \textit{sprints}, $c_p = \text{R\$}\,2.000$ por ponto e
$C = \text{R\$}\,480.000$.

A simplicidade da Equação~\ref{eq:previsao} esconde o problema econômico do
tema: cada termo herda a incerteza da estimativa que o produz. O
\textit{story point} é a métrica de tamanho mais usada em projetos ágeis
\cite{usman2014,fernandezdiego2020}, e sua atribuição costuma ser feita por
julgamento de especialistas, em especial com \textit{Planning Poker}
\cite{fernandezdiego2020}. Há indícios, porém, de que os pontos atribuídos por
humanos nem sempre refletem o tempo de desenvolvimento efetivamente gasto: em
37.440 histórias de 37 projetos de código aberto, \citeonline{tawosi2022}
encontram evidências de que o \textit{story point} pode não ser um bom
indicador de esforço. A
\textit{velocity}, por sua vez, não é constante, e por isso há trabalhos
dedicados a prevê-la \cite{hearty2009} e a antecipar se uma iteração em curso
vai entregar o que foi planejado \cite{choetkiertikul2018}. Por fim,
\citeonline{jorgensen2005} mostra que as avaliações de incerteza em custo de
software tendem ao excesso de confiança, ou seja, a intervalos mais estreitos
do que a realidade justifica. Erros nesses termos chegam ao orçamento, ao
cronograma e à decisão de investir.

A pesquisa recente tem tentado reduzir esse erro com aprendizado de máquina.
\citeonline{choetkiertikul2019} propõem o Deep-SE, que estima \textit{story
points} a partir do texto das \textit{issues} e superou as referências de
comparação em 23.313 \textit{issues} de 16 projetos de código aberto;
\citeonline{fu2023} seguem a mesma linha com um modelo baseado em
transformadores (GPT2SP). Entretanto, ao replicar o Deep-SE,
\citeonline{tawosi2023} constatam que ele supera as estimativas de mediana e o
TF-IDF com SVM em poucos casos com significância estatística. Esse debate
reforça a importância de comparar qualquer estimador com um \textit{baseline}
bem definido, prática que, para a estimativa de esforço em geral, ajuda a
reduzir a instabilidade das conclusões \cite{sarro2018}. Do lado gerencial, o uso das métricas ágeis em
escala também é desafiador: \citeonline{schull2023} mostram, em estudo de caso
numa grande montadora, a dificuldade de reportar progresso e desempenho de
equipes ágeis aos níveis de produto e portfólio. E, na ótica da economia da
informação, \citeonline{koch2008} ilustra, para projetos de código aberto, que
modelos de esforço baseados na participação de programadores estimam um esforço
sensivelmente menor do que modelos baseados em medidas de saída, como linhas de
código, isto é, a escolha da medida de esforço altera as conclusões sobre a
eficiência econômica do desenvolvimento.

\subsection{Objetivos}

O objetivo geral deste trabalho é analisar, com base na literatura, como as
estimativas baseadas em histórias de usuário (\textit{story points} e
\textit{velocity}) sustentam previsões de esforço, prazo e custo em projetos
ágeis, e quais são os limites dessas previsões. Constituem objetivos
específicos:
%
\begin{alineas}
  \item conceituar histórias de usuário, \textit{story points} e
        \textit{velocity};
  \item descrever como a produtividade da equipe converte pontos em prazo e em
        custo;
  \item examinar as evidências empíricas sobre a acurácia dessas estimativas;
  \item apresentar exemplos e aplicações, acadêmicos e de mercado; e
  \item discutir limitações e boas práticas para o uso financeiro das
        estimativas.
\end{alineas}

\subsection{Procedimentos metodológicos}

Trata-se de uma pesquisa bibliográfica, exploratória e qualitativa. Partiu-se
dos periódicos sugeridos no enunciado do seminário, a saber, \textit{IEEE
Transactions on Software Engineering} (TSE), \textit{ACM Transactions on
Software Engineering and Methodology} (TOSEM), \textit{Information Systems and
e-Business Management} (ISeBM) e \textit{Information Economics and Policy}
(IEP), complementados por revisões sistemáticas, anais de eventos e livros de
referência da área. Os dados bibliográficos (volume, número, páginas e DOI)
foram conferidos no Crossref. Nas buscas realizadas, os estudos dedicados a
\textit{story points} e \textit{velocity} concentraram-se no TSE; do TOSEM, do
ISeBM e do IEP foram aproveitados trabalhos adjacentes, sobre
\textit{baselines} de estimativa, reporte de métricas em organizações ágeis e
modelagem econômica de esforço, respectivamente. Essa distribuição desigual é
uma limitação do levantamento e não esgota a literatura disponível.

\subsection{Organização do trabalho}

Além desta introdução, o artigo está organizado em três seções. A
Seção~\ref{sec:marco} apresenta o marco teórico sobre histórias de usuário,
\textit{story points}, \textit{velocity} e sua relação com previsões
financeiras. A Seção~\ref{sec:exemplos} traz exemplos e aplicações, e a
Seção~\ref{sec:conclusao} reúne as conclusões. Seguem-se as referências
bibliográficas.
